\section{Exactly \texorpdfstring{$d-1$}{d-1} angular zeros on every upper semicircle}
\label{ado:sec:angular}
The signed transform gives, for $0<\rho\le1$ and $0<\theta<\pi$,
\begin{equation}\label{ado:eq:angulartransform}
 \Imn L(\rho e^{\ii\theta})
 =\rho\sin\theta\,J_\rho(\cos\theta),\qquad
 J_\rho(c)=\int_0^1
       \frac{\kappa(u)}{1-2\rho cu+\rho^2u^2}\dd u.
\end{equation}
A bound based only on simple sign-change arguments would leave open the
possibility of tangencies. We therefore prove a multiplicity-sensitive
upper bound before proving existence.

\subsection{An elementary confluent Cauchy-kernel bound}
\begin{lemma}[Angular variation bound with multiplicities]
\label{ado:lem:variation}
Suppose $\kappa\in L^1(0,1)$ is continuous, has exactly $m$ sign
changes at $r_1<\cdots<r_m$, and has a nonzero fixed sign on each
complementary interval. Then $J_\rho$ has at most $m$ zeros in
$(-1,1)$, counted with analytic multiplicity, for every $0<\rho\le1$.
\end{lemma}
\begin{proof}
The function
\begin{equation}\label{ado:eq:cauchymap}
 t(u)=\frac{1+\rho^2u^2}{2\rho u}
\end{equation}
is strictly decreasing on $(0,1)$, since
$t'(u)=(\rho^2u^2-1)/(2\rho u^2)<0$. Its values are at least one,
and it tends to infinity as $u\downarrow0$. Thus
\[
 J_\rho(c)=\int_0^1\frac{\kappa(u)}{2\rho u}
                         \frac1{t(u)-c}\dd u.
\]
The factor $\kappa(u)/u$ need not be integrable by itself; this causes
no difficulty because the displayed Cauchy factor is $O(u)$ at zero.
All the integrals below are interpreted with their complete kernels.

Suppose there are distinct zeros $c_1,\ldots,c_q\in(-1,1)$ with
positive multiplicities $\ell_1,\ldots,\ell_q$ summing to $m+1$.
If more zeros are available, retain any subcollection of total
multiplicity $m+1$. Form the proper rational function
\begin{equation}\label{ado:eq:variationrational}
 R(t)=\frac{\prod_{j=1}^m(t-t(r_j))}
             {\prod_{a=1}^q(t-c_a)^{\ell_a}}.
\end{equation}
Its partial fractions are linear combinations of
$(t-c_a)^{-h}$ for $1\le h\le\ell_a$. Differentiation of $J_\rho$
under the integral is valid locally for $c\in(-1,1)$; explicitly,
\[
 J_\rho^{(h-1)}(c)
 =(h-1)!\int_0^1\frac{\kappa(u)}{2\rho u}
                         \frac1{(t(u)-c)^h}\dd u.
\]
The selected zero multiplicities therefore imply
\begin{equation}\label{ado:eq:contradictoryintegral}
 \int_0^1\frac{\kappa(u)}{2\rho u}R(t(u))\dd u=0.
\end{equation}
This integral is absolutely convergent: near zero $R(t)=O(1/t)=O(u)$,
and near one its denominators stay positive and bounded away from zero.

But each factor $t(u)-t(r_j)$ changes sign exactly where $\kappa$
changes sign. The denominator of $R(t(u))$ is strictly positive, because
$t(u)\ge1>c_a$. Hence the integrand in
\eqref{ado:eq:contradictoryintegral} has one fixed nonzero sign except
at the $r_j$. Its integral cannot vanish. This contradiction proves the
bound and also excludes an identically zero $J_\rho$.
\end{proof}
The proof is a finite partial-fraction argument, including repeated poles.
No unproved strict-total-positivity assertion is being used.

\subsection{Positive compensation supplies the missing zeros}
For $0<\theta<\pi$ and $0\le u\le1$, put
\[
 \beta_u(\theta)=-\Arg(1-\rho u e^{\ii\theta}).
\]
For $0<u<1$ one has
\begin{equation}\label{ado:eq:betasector}
 0<\beta_u(\theta)<\beta_1(\theta)<\frac\pi2.
\end{equation}
The angle is strictly increasing in $u$, since differentiation with
respect to $u$ gives $\rho\sin\theta/|1-\rho u e^{\ii\theta}|^2>0$.
A positive mixture of the resolvents $(1-zu)^{-1}$ stays strictly inside
this sector. Consequently, with $H$ as in
\eqref{ado:eq:factorization},
\begin{equation}\label{ado:eq:Hsector}
 0<\Arg H(\rho e^{\ii\theta})<\beta_1(\theta).
\end{equation}
For example, the second inequality follows by rotating each resolvent
through $-\beta_1$ and observing that its imaginary part is strictly
negative for $0<u<1$. All rays lie in an angle of width less than
$\pi/2$, so no argument ambiguity arises.

\begin{theorem}[Complete angular count and signs]
\label{ado:thm:angular}
For \eqref{ado:eq:domain} and every $0<\rho\le1$,
$\Imn L(\rho e^{\ii\theta})$ has exactly $d-1$ simple zeros in
$0<\theta<\pi$. They have the signs and bounds stated in
Theorem~\ref{ado:thm:main}.
\end{theorem}
\begin{proof}
A continuous unwrapped argument of $L$ along the upper arc is
\begin{equation}\label{ado:eq:unwrappedphase}
 \Phi(\theta)=d\theta+
       \sum_{j=1}^{d-1}\beta_{r_j}(\theta)
       +\Arg H(\rho e^{\ii\theta}).
\end{equation}
The factorization proves this formula exactly; $\Phi$ is not reduced
modulo $2\pi$. The sector bounds imply
\begin{equation}\label{ado:eq:phasebounds}
 d\theta<\Phi(\theta)<d\bigl(\theta+\beta_1(\theta)\bigr)<d\pi.
\end{equation}
At the other end of the arc, $\Phi(\theta)\to d\pi$ as
$\theta\uparrow\pi$, since all resolvents approach positive real
values and $H(-\rho)>0$.

If $\rho<1$, then $\Phi(\theta)\to0$ as $\theta\downarrow0$.
The case $\rho=1$ requires more care: $L(1)$ may be singular, and
there is no reason to assign it argument zero. Instead, each fixed
$r_j<1$ gives $\beta_{r_j}(\theta)\to0$, while
\eqref{ado:eq:Hsector} gives
\[
 \limsup_{\theta\downarrow0}\Arg H(e^{\ii\theta})\le\pi/2.
\]
It follows that $0<\Phi(\theta)<\pi$ for sufficiently small positive
$\theta$. This weaker endpoint statement is all that is needed.

By continuity, the phase crosses each of $\pi,2\pi,\ldots,(d-1)\pi$.
These yield at least $d-1$ distinct zeros of the imaginary part. On the
other hand, Lemma~\ref{ado:lem:variation} with $m=d-1$, together
with \eqref{ado:eq:angulartransform}, gives at most $d-1$ zeros
counted with multiplicities. The change of variable $c=\cos\theta$
has nonzero derivative in the open arc, and the prefactor
$\rho\sin\theta$ is positive. Thus multiplicities are preserved.
Every zero is simple, and there is exactly one crossing of each listed
phase level. The levels must occur in their stated order, since a
reversal would force a repeated crossing of an earlier level.

The initial phase lies in $(0,\pi)$, so the imaginary part starts
positive. Simplicity forces alternating signs thereafter. At the
$k$th zero, the phase is $k\pi$ and the modulus is nonzero, so the real
value has sign $(-1)^k$.

For the location bounds, write $\alpha_k=k\pi/d$. At a zero,
\eqref{ado:eq:phasebounds} gives
\[
 \theta_k<\alpha_k<\theta_k+\beta_1(\theta_k).
\]
The function $h(\theta)=\theta+\beta_1(\theta)$ is strictly
increasing, because
\[
 h'(\theta)=\frac{1-\rho\cos\theta}
                  {1-2\rho\cos\theta+\rho^2}>0.
\]
For $\rho<1$, solving $h(\theta)=\alpha$ gives
$\theta=\alpha-\arcsin(\rho\sin\alpha)$. At $\rho=1$ the same
expression is the solution when $\alpha>\pi/2$; when
$\alpha\le\pi/2$ it is zero, and the strict lower bound follows
from $\theta_k>0$. This proves \eqref{ado:eq:mainangles}.
\end{proof}
The argument proves the needed crossings without asserting that
$\Phi$ is globally monotone. Such an assertion is unnecessary here.

\begin{corollary}[Analytic motion and no interior collisions]
\label{ado:cor:analyticbranches}
For fixed integer leading indices, each angular branch is real analytic
in $(\rho,b)\in(0,1)\times(0,\infty)$, and extends real analytically
across $\rho=1$ locally at each of its nontrivial upper-arc zeros.
The density roots $r_j$ are real analytic in $b>0$. No two of the
angular branches collide for $0<\rho\le1$.
\end{corollary}
\begin{proof}
On compact subsets of $\Ren b>0$, the seed $f_b$ is holomorphic as
an $L^1$-valued function; factors $|\log(-\log u)|^j$ arising from
parameter differentiation have an integrable common majorant on smaller
compact subsets. The bounded operators preserve this holomorphic
parameter dependence. The signed transform is locally holomorphic in
$z\in\Om$ and in $b$. Simplicity of its angular zeros permits the
real analytic implicit-function theorem. The unit-radius zeros remain
away from the cut, so the same local argument crosses $\rho=1$.
The density derivative formulas give joint real analyticity in $u,b$
in the open interval, and their simple roots obey the same theorem.
Distinctness follows from Theorem~\ref{ado:thm:angular}.
\end{proof}
The local extension across radius one does not imply that the global
count persists at all larger radii; Section~\ref{ado:sec:audit} gives
an explicit obstruction.

\subsection{Strict radial motion of every angular branch}
The compensation yields more than a zero count. It also turns a
monotonicity property of ordinary positive Stieltjes transforms into
strict motion of all the strict-depth angular zeros.

\begin{lemma}[Radial argument monotonicity for positive transforms]
\label{ado:lem:radialPick}
Let $\nu$ be a nonzero finite positive measure on $[0,1]$ and
$Q(z)=\int(1-uz)^{-1}\dd\nu(u)$. For $0<\theta<\pi$ and
$0<\rho\le1$,
\begin{equation}\label{ado:eq:radialPick}
 \frac{\partial}{\partial\rho}\Arg Q(\rho e^{\ii\theta})\ge0.
\end{equation}
The argument is chosen continuously from $Q(0)>0$.
\end{lemma}
\begin{proof}
First take a finite measure with positive weights at distinct points
$0<u_1<\cdots<u_n<1$. The rational function
$\sum_j w_j/(t-u_j)$ is strictly decreasing between consecutive poles,
with limits $+\infty$ and $-\infty$ at the two ends. Its numerator
therefore has exactly one root $v_j\in(u_j,u_{j+1})$ in each such
interval. Factoring the numerator gives
\[
 Q(z)=Q(0)\frac{\prod_{j=1}^{n-1}(1-v_jz)}
                      {\prod_{j=1}^n(1-u_jz)}.
\]
Consequently the radial derivative of its argument is
\[
 \sin\theta\left\{f(u_1)+
          \sum_{j=1}^{n-1}\bigl(f(u_{j+1})-f(v_j)\bigr)\right\},
 \qquad f(t)=\frac{t}{1-2\rho t\cos\theta+\rho^2t^2}.
\]
Every summand is nonnegative, because
\[
 f'(t)=\frac{1-\rho^2t^2}
            {(1-2\rho t\cos\theta+\rho^2t^2)^2}\ge0
 \qquad(0\le t\le1).
\]

Approximate an arbitrary finite positive $\nu$ weakly by such finite
measures, with their masses converging to $\nu([0,1])$. For each
compact subset of $\Om$, the integrands defining the transform and
its first complex derivative are bounded and continuous uniformly in
$u\in[0,1]$. The transforms and derivatives converge locally uniformly.
Also $Q(\rho e^{\ii\theta})\ne0$: it has positive imaginary part
unless the measure is supported at zero, in which case it is a positive
constant. Passing to the limit in
$\Imn(e^{\ii\theta}Q'(z)/Q(z))$ proves the nonnegative derivative.
The same argument works at $\rho=1$ because $\theta$ is strictly
between zero and $\pi$.
\end{proof}

\begin{theorem}[Strict inward angular motion]
\label{ado:thm:radialmotion}
For $d\ge2$ and \eqref{ado:eq:domain}, every angular branch satisfies
\begin{equation}\label{ado:eq:strictmotion}
 \frac{\dd\theta_k}{\dd\rho}<0,
 \qquad 0<\rho\le1,\quad 1\le k\le d-1.
\end{equation}
The derivative at one is the one-sided derivative, or equivalently the
derivative of the local continuation. The same result holds for the
positive-seed and shifted-terminal families of
Theorem~\ref{ado:thm:seedextension}.
\end{theorem}
\begin{proof}
At fixed $\theta$, differentiate the unwrapped phase
\eqref{ado:eq:unwrappedphase} with respect to the radius:
\[
 \Phi_\rho
 =\sum_{j=1}^{d-1}
      \frac{r_j\sin\theta}{1-2\rho r_j\cos\theta+\rho^2r_j^2}
   +\partial_\rho\Arg H(\rho e^{\ii\theta})>0.
\]
The first sum is strictly positive since $d\ge2$ and $r_j>0$; the
last term is nonnegative by Lemma~\ref{ado:lem:radialPick}.
At the unique simple crossing of the phase level $k\pi$, one has
$\Phi_\theta>0$. Indeed the phase starts below that level, ends
above it, and crosses it exactly once; simplicity excludes a zero
phase derivative there. Implicit differentiation of
$\Phi(\theta_k(\rho),\rho)=k\pi$ now gives
\[
 \theta_k'(\rho)=-\frac{\Phi_\rho}{\Phi_\theta}<0.
\]
All ingredients also hold for a general positive seed.
\end{proof}
For depth two, this does \emph{not} settle the manuscript's conjecture
about $\eta(\rho)=\cos\theta(\rho)/\rho$: the derivative of that
quotient has a second term, and its sign is not implied by
$\theta'(\rho)<0$. The two kinds of radial monotonicity must remain
separate in the integration ledger.
